\section{Background}
\label{sec:cohear_background}

This section lays the technical foundation for CoHear, focusing on acoustic sensing techniques used for conversation modeling.
First, it discusses sound localization, which estimates speaker positions from earphone microphone arrays and provides information for constructing a conversation map.
It then reviews wireless acoustic sensor networks (WASNs), which support collaborative sensing across multiple devices.


\subsection{Sound Source Localization}
Sound source localization is a core problem in multi-channel audio and enables spatial acoustic sensing.
Consider a microphone array consisting of $N$ microphones located at known positions $\mathbf{p}_i = (x_i, y_i)$ for $i = 1, 2, \ldots, N$. Let the position of the sound source be denoted by $\mathbf{s} = (x_s, y_s)$.
The time it takes for the sound to reach the microphone $i$ can be expressed as:
\(t_{s, i} = \frac{d_{s, i}}{c}, d_{s, i} = \sqrt{(x_s - x_i)^2 + (y_s - y_i)^2}\)
where $d_{s, i}$ is the distance from sound source $s$ to microphone $i$, and $c$ is the speed of sound in air. Since the source emission time is unknown to the receiver, $t_{s, i}$ cannot be estimated directly. Instead, the system estimates the time difference of arrival (TDoA), $t_{s, i} - t_{s, j}$, where $i, j \in (1, N)$. A common method for estimating TDoA from in-the-wild recordings is the generalized cross-correlation (GCC)~\cite{knapp1976generalized}, which is computed for each microphone pair. The GCC peak estimates the TDoA, which can then be used to calculate the direction of arrival (DoA) as \(DoA = arctan((x_s - x_c)/(y_s - y_c))\), where $x_c$ is the center of the microphone array. This analysis shows why localization performance depends strongly on the number of microphones: more microphones provide more microphone pairs and therefore more TDoA constraints.

Instead of relying on analytical methods, sound source localization can be performed using a deep neural network, as explored in \cite{wang2022nerc}. This deep learning-based approach generally consists of two stages: feature extraction and location estimation.
In the feature extraction stage, the audio signal is typically converted into a spectrogram representation with dimensions: \(C\times T \times F\), where \(C\) refers to channels, \(T\) refers to time, and \(F\) refers to frequency. These features can be further processed using techniques such as log-melspectrogram conversion, which is applied to each channel individually, or multi-channel methods like GCC \cite{knapp1976generalized} and SALSA-Lite \cite{nguyen2022salsa}.
In the location estimation stage, the neural network is trained to output the sound source location, which also supports a dynamic number of simultaneous sound sources. For instance, architectures employing multi-track \cite{cao2021improved} or multi-ACCDOA \cite{shimada2022multi} output formats are designed specifically to localize overlapping sound events.

Deep sound localization can also use the spatial information embedded in binaural recordings, including reflections from the human head and body. Specifically, the head-related transfer function (HRTF) can be modeled as:
\(H_{L}(\theta, \phi, f, d) = \frac{Y_{L}(f)}{X(f)}, H_{R}(\theta, \phi, f, d) = \frac{Y_{R}(f)}{X(f)},\) where $H_{L}$ indicates the frequency response for left ear and $H_{R}$ indicates the frequency response for right ear. 
By using HRTF cues with deep learning, binaural localization can achieve higher accuracy than a linear two-microphone array. For instance, this approach can partially resolve front--back confusion, a problem that is particularly challenging for conventional two-microphone arrays.


\subsection{Wireless Acoustic Sensor Network}
We are studying a wireless acoustic sensor network where a group of sensor nodes is randomly distributed in a reverberant environment \cite{bertrand2011applications}. We assume that the internal geometric configuration of each node's microphone array is known and that all microphones within the array are sampled at the same time, which we believe is a reasonable assumption following \cite{gburrek2021geometry}. Additionally, we assume that there is coarse time synchronization among the clocks of the sensor nodes, enabling us to aggregate the data globally. By default, we assume a two-dimensional space, but it can be easily extended to a three-dimensional space.

Since a sensor node cannot determine its own position or orientation within the global coordinate system, all observations are made in the local coordinate system of each node. 
Each sensor node \( l \) (where \( l = 1, \ldots, L \)) computes acoustic features (e.g., DoA estimates \( \phi_k^l \) and distance \( d_k^l \)) to the acoustic source \( k \), all reference to the node's local coordinate system. Consequently, we gather a total of \( K \times L \) estimates as the inputs.

From the global view of the WASN, it is composed of \( L \) sensor nodes each equipped with a microphone array located at positions \( n_l = [n_{l,x}, n_{l,y}]^T \) and oriented at angles \( \theta_l \) (where \( l = 1, 2, \ldots, L \)) relative to a global coordinate system defined by the x and y axes. There are \( K \) acoustic sources located at positions \( s_k = [s_{k,x}, s_{k,y}]^T \) (where \( k = 1, 2, \ldots, K \)). Note that all the above global notations are unknown and will be estimated through a geometry calibration process based on the observed acoustic signals from sensor nodes.



